问题一先要验证"鱼类资源量增加"，还要解释"鱼类增长的驱动机制"和"基础资源的承载压力变化"。若把四大家鱼并成一个消费者功能群，再将水草、浮游植物、浮游动物和底栖饵料合成单一资源总量，模型虽仍能贴合2025年资源量恢复至1.8倍的观测值，但无法判断哪一层资源先承压，也难以解释"水下荒漠"这类底层生态退化。因此，正式模型必须保留四类基础资源及其与四大家鱼的基本食性对应关系。
